\documentclass[10.5pt]{article}
\usepackage[a4paper,margin=2cm]{geometry}
\usepackage[T1]{fontenc}
\usepackage{lmodern}
\usepackage{microtype}
\usepackage{enumitem}
\usepackage{xcolor}
\usepackage[hidelinks]{hyperref}
\usepackage{titlesec}
\definecolor{accent}{HTML}{1F4E79}
\titleformat{\section}{\color{accent}\large\bfseries}{}{0pt}{}
\titlespacing*{\section}{0pt}{10pt}{4pt}
\setlist{nosep,leftmargin=1.2em}
\setlength{\parindent}{0pt}
\setlength{\parskip}{4pt}
\pagestyle{empty}

\begin{document}

{\color{accent}\LARGE\bfseries Polygenic Liability to Bipolar Disorder and the Plasma Proteome}\\[3pt]
{\large Research Project Overview}

\vspace{6pt}
\textbf{Yunqi (Vicky) Liu} \quad MSc Medical Genetics and Genomics, University of Glasgow\\
MSc dissertation project, June--August 2025 \quad Supervisor: Dr Rona Strawbridge\\
\textbf{Status:} manuscript in preparation\\
Code: \url{https://github.com/bleulievre/bipolar-prs-proteomics} \quad Contact: \href{mailto:vickyliu1998424@gmail.com}{vickyliu1998424@gmail.com}

\section{Background}
Bipolar disorder is highly heritable, yet its underlying mechanisms remain unclear and clinically useful blood
biomarkers are still an open research question. Polygenic risk scores (PRS) summarize genetic liability across
the genome and can be calculated for individuals with or without a diagnosis, so they allow liability to be related to
biological measurements in a large population cohort. Large-scale plasma proteomics in UK Biobank now makes it
possible to ask whether this liability is reflected in circulating protein levels.

\section{Research question}
Is polygenic liability to bipolar disorder associated with the levels of circulating plasma proteins in a
population cohort, and do any such associations remain after adjustment for demographic, technical,
lifestyle, cardiometabolic and psychiatric factors?

\section{Data}
\begin{itemize}
  \item UK Biobank participants with genotype and Olink plasma proteomic data (49,239 participants in the
        dataset supplied by the research group; about 730 protein assays).
  \item Bipolar disorder PRS: LDpred scores from published Psychiatric Genomics Consortium summary statistics,
        provided by the research group.
  \item Covariates: age, sex, genetic principal components, genotyping array, deprivation, BMI, waist-to-hip
        ratio, smoking, alcohol, blood pressure and medication, diabetes, cardiovascular history, and
        psychiatric history and medication.
\end{itemize}

\section{Analysis design}
\begin{enumerate}
  \item \textbf{Quality control.} Profiled missingness for every assay and set an exclusion threshold for
        poorly measured proteins.
  \item \textbf{Screening.} Spearman rank correlation between the PRS and each protein, with a
        Bonferroni-based threshold for the number of assays tested.
  \item \textbf{Sequential multivariable regression.} Linear models for the screened proteins, adding covariate
        groups step by step (basic, then lifestyle, cardiometabolic and medication covariates), with the
        Bonferroni threshold reset at each step.
  \item \textbf{Sensitivity analysis.} Repeated the final model after excluding participants with any lifetime
        mental illness or psychiatric medication use, to assess whether the associations were robust to
        excluding them.
  \item \textbf{Biological interpretation.} Gene Ontology enrichment (clusterProfiler) and literature-based
        functional annotation of the proteins that remained associated.
\end{enumerate}
A small set of candidate proteins remained associated after full adjustment and the sensitivity analysis.
The specific proteins and estimates are withheld until the manuscript is published.

\section{My role}
\begin{itemize}
  \item Cleaned and merged the UK Biobank extract and the psychiatric-history file, coded covariates, and set
        the missing-data threshold.
  \item Implemented the staged analysis in R under Dr Strawbridge's supervision, and interpreted the results
        against published genetic and clinical evidence.
  \item Wrote the dissertation, with a log of investigation recording what was done, why, and the
        alternatives considered.
\end{itemize}

\section{Skills demonstrated}
\begin{itemize}
  \item \textbf{Statistics:} rank correlation, multiple-testing correction, multivariable linear regression,
        sequential covariate adjustment, sensitivity analysis.
  \item \textbf{Biobank data:} large-scale phenotype and proteomic data handling, missingness profiling,
        covariate harmonization, linking auxiliary phenotype files.
  \item \textbf{Programming:} R (dplyr, tidyr, ggplot2, clusterProfiler, org.Hs.eg.db).
  \item \textbf{Interpretation:} pathway enrichment and protein-level literature review in the context of
        psychiatric genetics.
\end{itemize}

\section{Current status}
A manuscript is being prepared with Dr Strawbridge, and the analyses will be re-run within the approved
UK Biobank project before submission. Separately, after the MSc, the dissertation analysis was re-implemented
with AI coding assistance as a tested, configuration-driven pipeline that follows the original code step by step
and runs end-to-end on simulated data.

\section{What is available}
\begin{itemize}
  \item \textbf{Public:} a public copy of the dissertation R script (comments translated, protein names
        removed) and the re-implemented pipeline with simulated data and tests, at the repository above. No
        participant-level data, protein identities, effect estimates or findings are included.
  \item \textbf{On request:} further materials, including the dissertation, can be shared with the
        supervisor's agreement.
\end{itemize}

\end{document}
